\documentclass[12pt]{article}
\usepackage[utf8]{inputenc}
\usepackage[T1]{fontenc}
\usepackage[margin=1in]{geometry}
\usepackage{setspace}
\usepackage{xcolor}
\usepackage[colorlinks=true, urlcolor=blue!60!black, linkcolor=blue!60!black]{hyperref}
\setlength{\parindent}{2em}
\setlength{\parskip}{1.25em}
\renewcommand{\baselinestretch}{1.0}
\begin{document}
\begin{center}
{\LARGE \textit{United States v. Lopez}, 514 U.S. 549 (1995)\par}
\vspace{0.8em}
{\large Case Brief\par}
\end{center}
\vspace{1em}

\section{Memory Jogger}
For the first time in nearly six decades, the Supreme Court struck down a federal statute as exceeding Congress's Commerce Clause power, holding that possessing a gun in a local school zone is neither an economic activity nor tied by a jurisdictional element to interstate commerce (\href{https://www.courtlistener.com/opinion/117927/united-states-v-lopez/}{CourtListener}; \href{https://www.law.cornell.edu/supct/html/93-1260.ZO.html}{Cornell LII}).

\section{Detailed Case Facts}
On March 10, 1992, Alfonso Lopez, Jr., a twelfth-grade student at Edison High School in San Antonio, Texas, arrived at school carrying a concealed .38-caliber handgun and five bullets (\href{https://www.law.cornell.edu/supct/html/93-1260.ZO.html}{Majority Opinion}). Acting on an anonymous tip, school authorities confronted Lopez, who admitted he was carrying the firearm. He was arrested and charged under Texas law with firearm possession on school premises (Tex. Penal Code Ann. \S 46.03(a)(1) (Supp. 1994)).

The following day, state prosecutors dismissed the state charges after federal agents charged Lopez by criminal complaint with violating the Gun-Free School Zones Act of 1990 (GFSZA), 18 U.S.C. \S 922(q)(1)(A) (1988 ed., Supp. V) (\href{https://www.law.cornell.edu/supct/html/93-1260.ZO.html}{Majority Opinion}; \href{https://www.oyez.org/cases/1994/93-1260}{Oyez}). Section 922(q)(1)(A) made it a federal offense ``for any individual knowingly to possess a firearm at a place that the individual knows, or has reasonable cause to believe, is a school zone,'' with ``school zone'' defined in \S 921(a)(25) as in or on the grounds of, or within 1,000 feet of, a public, parochial, or private school. As originally enacted in 1990, neither the statute nor its legislative history contained express congressional findings linking school-zone gun possession to interstate commerce, nor did the statute require proof that the firearm had moved in or otherwise affected interstate commerce (\href{https://www.law.cornell.edu/supct/html/93-1260.ZO.html}{Majority Opinion}).

\section{Procedural History}
A federal grand jury indicted Lopez on one count of knowing possession of a firearm in a school zone in violation of 18 U.S.C. \S 922(q). Lopez moved to dismiss the indictment, arguing that \S 922(q) was unconstitutional because it exceeded Congress's power to legislate control over public schools (\href{https://www.law.cornell.edu/supct/html/93-1260.ZO.html}{Majority Opinion}). The United States District Court for the Western District of Texas denied the motion, concluding that \S 922(q) was a constitutional exercise of Congress's Commerce Clause authority because the ``business'' of elementary, middle, and high schools affects interstate commerce. Lopez waived his right to a jury trial; following a bench trial, the District Court found him guilty and sentenced him to six months' imprisonment and two years' supervised release (\href{https://www.oyez.org/cases/1994/93-1260}{Oyez}).

On appeal, the United States Court of Appeals for the Fifth Circuit reversed Lopez's conviction, 2 F.3d 1342 (5th Cir. 1993), holding that in light of insufficient congressional findings and legislative history, \S 922(q) exceeded Congress's authority under the Commerce Clause (\href{https://www.law.cornell.edu/supremecourt/text/93-1260}{Syllabus}). The Supreme Court granted certiorari, 511 U.S. 1029 (1994), heard oral argument on November 8, 1994, and affirmed the Fifth Circuit on April 26, 1995 (\href{https://www.oyez.org/cases/1994/93-1260}{Oyez}).

\section{Judicial Votes}
\textbf{Vote:} 5--4 affirming the Fifth Circuit (\href{https://www.law.cornell.edu/supremecourt/text/93-1260}{Syllabus}; \href{https://www.courtlistener.com/opinion/117927/united-states-v-lopez/}{CourtListener}).
\begin{itemize}
\item \textbf{Majority Opinion:} Chief Justice Rehnquist, joined by Justices O'Connor, Scalia, Kennedy, and Thomas (\href{https://www.law.cornell.edu/supct/html/93-1260.ZO.html}{Majority Opinion}).
\item \textbf{Concurring Opinions:} Justice Kennedy, joined by Justice O'Connor (\href{https://www.law.cornell.edu/supct/html/93-1260.ZC.html}{Kennedy Concurrence}); Justice Thomas (\href{https://www.law.cornell.edu/supct/html/93-1260.ZC1.html}{Thomas Concurrence}).
\item \textbf{Dissenting Opinions:} Justice Stevens (\href{https://www.law.cornell.edu/supct/html/93-1260.ZD.html}{Stevens Dissent}); Justice Souter (\href{https://www.law.cornell.edu/supct/html/93-1260.ZD1.html}{Souter Dissent}); Justice Breyer, joined by Justices Stevens, Souter, and Ginsburg (\href{https://www.law.cornell.edu/supct/html/93-1260.ZD2.html}{Breyer Dissent}).
\end{itemize}

\section{Holding}
\textbf{Judgment:} Affirmed. The Gun-Free School Zones Act of 1990, 18 U.S.C. \S 922(q)(1)(A), exceeds Congress's authority under the Commerce Clause, U.S. Const., Art. I, \S 8, cl. 3 (\href{https://www.law.cornell.edu/supct/html/93-1260.ZO.html}{Majority Opinion}).

\textbf{Rule Announced:} Congress may regulate three broad categories of activity under the Commerce Clause: (1) the use of the channels of interstate commerce; (2) the instrumentalities of interstate commerce, or persons or things in interstate commerce, even though the threat may come only from intrastate activities; and (3) intrastate activities that ``substantially affect'' interstate commerce (\href{https://www.law.cornell.edu/supct/html/93-1260.ZO.html}{Majority Opinion}). Under the third category, a criminal statute regulating purely local, noneconomic conduct (gun possession in a school zone) that is not an essential part of a larger regulation of economic activity, contains no express jurisdictional element tying the regulated conduct to interstate commerce, and relies on an attenuated ``inference upon inference'' causal chain to national economic effects cannot be sustained under the Commerce Clause.

\textbf{Holding vs. Dicta:} The narrow holding is that 18 U.S.C. \S 922(q)(1)(A), as enacted in 1990 without a jurisdictional hook and targeting noncommercial gun possession, exceeds the Commerce Clause. The Court's remarks that Congress is normally not required to make formal findings (and its observation in footnote 4 about Congress's 1994 post-prosecution findings, which the Government did not strictly rely upon) are dicta, as are its suggestions about hypothetical federal regulation of family law or school curricula (\href{https://www.law.cornell.edu/supct/html/93-1260.ZO.html}{Majority Opinion}). Likewise, the Court did not definitively hold that a jurisdictional element alone would automatically validate \S 922(q), though Congress subsequently amended \S 922(q) in 1996 to require that the firearm ``has moved in or that otherwise affects interstate or foreign commerce'' (\href{https://www.govregs.com/uscode/expand/title18_partI_chapter44_section922}{18 U.S.C. \S 922(q)}).

\section{Analysis of Opinions}

\subsection{Majority Opinion (Chief Justice Rehnquist)}
Chief Justice Rehnquist began from ``first principles'': the Constitution establishes a federal government of enumerated powers (citing James Madison in \textit{The Federalist No. 45}), and the division of authority between federal and state governments protects individual liberty (\href{https://www.law.cornell.edu/supct/html/93-1260.ZO.html}{Majority Opinion}). Surveying Commerce Clause history from \textit{Gibbons v. Ogden}, 9 Wheat. 1 (1824), through the pre-1937 direct/indirect era and the post-1937 expansion in \textit{NLRB v. Jones \& Laughlin Steel Corp.}, 301 U.S. 1 (1937), \textit{United States v. Darby}, 312 U.S. 100 (1941), and \textit{Wickard v. Filburn}, 317 U.S. 111 (1942), Rehnquist emphasized that even modern precedents recognize judicially enforceable outer limits.

Synthesizing prior case law, Rehnquist articulated the now-canonical three-category framework: (1) channels of interstate commerce; (2) instrumentalities of, or persons or things in, interstate commerce; and (3) activities having a ``substantial relation'' to interstate commerce---clarifying that the proper test under category three is whether the activity ``substantially affects'' (not merely ``affects'') interstate commerce (\href{https://www.law.cornell.edu/supct/html/93-1260.ZO.html}{Majority Opinion}).

Quickly disposing of categories one and two, Rehnquist evaluated \S 922(q) under category three and identified three fatal defects. First, \S 922(q) was a criminal statute having nothing to do with ``commerce'' or any economic enterprise, nor was it an essential part of a larger regulation of economic activity (distinguishing \textit{Wickard}, which regulated wheat production in a national market scheme). Second, \S 922(q) lacked an express jurisdictional element ensuring through case-by-case inquiry that the firearm possession affected interstate commerce (contrasting \textit{United States v. Bass}, 404 U.S. 336 (1971)). Third, while formal congressional findings are not required, their absence meant no legislative judgment was available to help the Court discern a substantial effect ``not visible to the naked eye.'' Rejecting the Government's ``costs of crime'' and ``national productivity'' theories, Rehnquist reasoned that accepting such attenuated causal chains would allow Congress to regulate all violent crime, family law, and local school curricula, obliterating the constitutional distinction between ``what is truly national and what is truly local'' and converting the Commerce Clause into a general federal police power (\href{https://www.law.cornell.edu/supct/html/93-1260.ZO.html}{Majority Opinion}).

\subsection{Concurring Opinion (Justice Kennedy, joined by Justice O'Connor)}
Justice Kennedy wrote separately to emphasize that the Court's holding was ``necessary though limited'' (\href{https://www.law.cornell.edu/supct/html/93-1260.ZC.html}{Kennedy Concurrence}). Tracing the Court's historical struggle to abandon formalistic pre-1937 categories, Kennedy stressed that \textit{stare decisis} counsels strongly against unsettling post-New Deal doctrine permitting Congress to regulate commercial transactions in a unified national market.

Nevertheless, federalism is the Framers' unique structural contribution to political science, protecting liberty through ``two distinct and discernable lines of political accountability'' (\href{https://www.law.cornell.edu/supct/html/93-1260.ZC.html}{Kennedy Concurrence}). When Congress regulates neither commercial actors nor commercial conduct and intrudes into an area of traditional state concern---here, education and local criminal law enforcement, where over 40 States already legislated and served as ``laboratories for experimentation''---the Court has a duty to intervene so that federal-state boundaries do not blur and political accountability does not become illusory.

\subsection{Concurring Opinion (Justice Thomas)}
Justice Thomas joined the majority opinion but urged the Court in a future case to reconsider the entire ``substantial effects'' test as a departure from the original public meaning of the Commerce Clause (\href{https://www.law.cornell.edu/supct/html/93-1260.ZC1.html}{Thomas Concurrence}). Citing Founding-era dictionaries (Samuel Johnson, Nathan Bailey, Thomas Sheridan), \textit{The Federalist}, and ratification debates, Thomas argued that at the Founding ``commerce'' meant selling, buying, bartering, and transporting for those purposes---distinct from productive activities such as agriculture and manufacturing.

Thomas contended that the substantial effects test and its aggregation principle render much of Article I, \S 8 (such as the bankruptcy, coinage, post office, patent, and piracy powers) superfluous surplusage and turn the Tenth Amendment on its head (\href{https://www.law.cornell.edu/supct/html/93-1260.ZC1.html}{Thomas Concurrence}). He also argued that \textit{Gibbons v. Ogden} had been misread by New Deal cases and that the ``wrong turn'' in Commerce Clause jurisprudence occurred not in the pre-1937 era, but in the 1930s departure from 150 years of precedent.

\subsection{Dissenting Opinion (Justice Stevens)}
In a brief solo dissent, Justice Stevens endorsed both Justice Breyer's empirical Commerce Clause analysis and Justice Souter's warning about a return to pre-1937 substantive due process (\href{https://www.law.cornell.edu/supct/html/93-1260.ZD.html}{Stevens Dissent}). Stevens added that firearms are both articles of commerce and articles that can be used to restrain commerce; because their possession is the direct or indirect consequence of commercial activity, Congress's power to regulate commerce in firearms includes prohibiting their possession in particular markets, including the substantial market for handguns among school-age children.

\subsection{Dissenting Opinion (Justice Souter)}
Justice Souter warned that the majority's heightened scrutiny of noncommercial regulation resurrected the discredited pre-1937 formalism of \textit{Lochner v. New York}, 198 U.S. 45 (1905), and \textit{Carter v. Carter Coal Co.}, 298 U.S. 238 (1936) (\href{https://www.law.cornell.edu/supct/html/93-1260.ZD1.html}{Souter Dissent}). Under modern rational basis review, courts ask only whether Congress had a rational basis for concluding that a regulated activity substantially affects interstate commerce and whether the chosen means are reasonably adapted to that end.

Souter rejected the majority's commercial/noncommercial distinction as porous and unworkable, rejected any notion that Congress's plenary commerce power diminishes when touching areas of traditional state concern, and cautioned against making express congressional findings dispositive, which would covertly substitute judicial review of legislative wisdom or deliberateness for deferential rationality review (\href{https://www.law.cornell.edu/supct/html/93-1260.ZD1.html}{Souter Dissent}).

\subsection{Dissenting Opinion (Justice Breyer, joined by Justices Stevens, Souter, and Ginsburg)}
Justice Breyer delivered the principal dissent, applying three settled principles: (1) Congress may regulate local activities that significantly affect interstate commerce; (2) courts evaluate the cumulative, aggregate effect of all similar instances; and (3) courts defer to Congress under rational basis review (\href{https://www.law.cornell.edu/supct/html/93-1260.ZD2.html}{Breyer Dissent}).

Marshaling an extensive appendix of congressional hearings, federal reports, and economic studies, Breyer argued that Congress had an ample rational basis for linking school-zone gun violence to classroom disruption, reduced educational attainment, diminished workforce skills in an integrated global economy, and business siting decisions (\href{https://www.law.cornell.edu/supct/html/93-1260.ZD2.html}{Breyer Dissent}). He contended that the majority's holding created three serious legal problems: it conflicted with precedents like \textit{Perez v. United States}, 402 U.S. 146 (1971), \textit{Katzenbach v. McClung}, 379 U.S. 294 (1964), and \textit{Wickard}; it relied on an unworkable commercial/noncommercial line (noting public schools spent \$230 billion in 1990); and it injected legal uncertainty into over 100 federal statutes using ``affecting commerce'' language or lacking jurisdictional hooks.

\section{Examples: Future Applications}
\begin{enumerate}
\item \textbf{Same-Side (Invalid under \textit{Lopez}): Federal Ban on Simple Possession of Knuckles or Switchblades in Public Parks.} Suppose Congress enacts a standalone criminal statute barring possession of brass knuckles within 500 feet of any municipal park, with no jurisdictional element and no link to a broader commercial regulatory scheme. Like \S 922(q), mere local possession of a weapon in a noncommercial setting is not economic activity, and justifying it through general ``costs of crime'' or reduced tourism would require piling inference upon inference.
\item \textbf{Same-Side (Invalid under \textit{Lopez}): Federal Criminalization of Intrastate Truancy.} Suppose Congress makes it a federal misdemeanor for a high school student to miss more than ten unexcused days of public school, citing national workforce productivity. Under \textit{Lopez}, truancy is purely noneconomic conduct in an area of traditional state control (education), and the ``national productivity'' rationale was expressly rejected as lacking any limiting principle.
\item \textbf{Opposite-Side (Valid under \textit{Lopez}): Federal Ban on Intrastate Cultivation of a Controlled Substance Within a Comprehensive Market Regulation.} If Congress prohibits purely local, noncommercial home cultivation of marijuana as part of a comprehensive federal regulatory scheme controlling the interstate market for controlled substances, the statute falls on the opposite side of \textit{Lopez}: it regulates an ``essential part of a larger regulation of economic activity, in which the regulatory scheme could be undercut unless the intrastate activity were regulated'' (as the Court later held in \href{https://www.law.cornell.edu/supremecourt/text/03-1454}{\textit{Gonzales v. Raich}, 545 U.S. 1 (2005)}).
\item \textbf{Opposite-Side (Valid under \textit{Lopez}): Federal Prohibition on Shipping Unregistered Firearms Across State Lines or Destroying Interstate Rail Bridges.} A federal law prohibiting interstate shipment of unserialized firearms directly regulates the \textit{channels} of interstate commerce (category one), and a statute criminalizing sabotage of interstate freight rail tracks protects the \textit{instrumentalities} of interstate commerce (category two), both of which \textit{Lopez} expressly reaffirmed as squarely within Congress's power (\href{https://www.law.cornell.edu/supct/html/93-1260.ZO.html}{Majority Opinion}).
\item \textbf{Fence-Sitter (Close Call): Reenacted Gun-Free School Zones Act with a Minimal Jurisdictional Hook.} Suppose Congress amends \S 922(q)---as it did in 1996 (\href{https://www.govregs.com/uscode/expand/title18_partI_chapter44_section922}{18 U.S.C. \S 922(q)(2)(A)})---to criminalize possessing in a school zone ``a firearm that has moved in or that otherwise affects interstate or foreign commerce.'' On one hand, adding an express jurisdictional element addresses Rehnquist's second defect and aligns with \textit{Scarborough v. United States}, 431 U.S. 563 (1977), under which prior interstate movement of a gun suffices. On the other hand, critics argue that if a firearm's crossing of a state line decades earlier automatically authorizes federal prosecution of local possession, the jurisdictional hook becomes a formalistic loophole that undermines \textit{Lopez}'s substantive federalism limit.
\end{enumerate}

\section{Critique}
\textit{United States v. Lopez} sparked an intense scholarly debate over federalism, judicial review, and constitutional methodology. From an \textbf{originalist and textualist perspective}, scholars such as Randy Barnett have defended \textit{Lopez} as a legitimate, non-activist restoration of the Constitution's enumerated-powers structure, arguing that original public meaning supports an even narrower reading of ``commerce'' along the lines of Justice Thomas's concurrence---confined to trade and transportation rather than all gainful economic activity (\href{https://scholarship.law.georgetown.edu/context/facpub/article/1504/viewcontent/barnett_commerce_clause.pdf}{Barnett, 68 U. Chi. L. Rev. 101 (2001)}; \href{https://scholarship.law.georgetown.edu/cgi/viewcontent.cgi?article=1856&context=facpub}{Barnett, 72 U. Colo. L. Rev. 1275 (2002)}). At the same time, some originalists critique Rehnquist's majority opinion as an uneasy half-measure: by preserving \textit{Wickard} and the ``substantial effects'' framework while drawing an amorphous line between ``economic'' and ``noneconomic'' activity, the Court left the doctrine vulnerable to later expansion in \href{https://www.law.cornell.edu/supremecourt/text/03-1454}{\textit{Gonzales v. Raich}} (2005).

From a \textbf{progressive and doctrinal perspective}, critics such as Erwin Chemerinsky and commentators echoing Justices Souter and Breyer contend that \textit{Lopez} revived arbitrary pre-1937 formalisms without a principled or predictable standard for distinguishing ``economic'' from ``noneconomic'' conduct (\href{http://www.stanfordlawreview.org/wp-content/uploads/sites/3/2010/04/chemerinsky.pdf}{Chemerinsky, 58 Stan. L. Rev. 1763 (2006)}; \href{https://scholarlycommons.law.case.edu/cgi/viewcontent.cgi?article=4685&context=caselrev}{Case W. Res. L. Rev.}). Progressive scholars also challenge the majority's assumption that judicial enforcement of federalism enhances liberty or political accountability, arguing instead that national problems in an integrated economy require flexible national democratic solutions and that the political process---not unelected judges---should police the federal-state balance.

Underlying these methodological disagreements is a deep clash of \textbf{competing constitutional values}: the structural commitment to limited, enumerated federal powers, local experimentation, and clear lines of dual-sovereign accountability versus democratic deference to Congress's superior empirical capacity to address nationwide social and economic harms in a deeply interconnected modern society.

\section{Key Quotations}
\begin{quote}
``Section 922(q) is a criminal statute that by its terms has nothing to do with `commerce' or any sort of economic enterprise, however broadly one might define those terms. Section 922(q) is not an essential part of a larger regulation of economic activity, in which the regulatory scheme could be undercut unless the intrastate activity were regulated.''
\end{quote}
\noindent --- Chief Justice Rehnquist, Opinion of the Court (\href{https://www.law.cornell.edu/supct/html/93-1260.ZO.html}{Cornell LII}; \href{https://www.courtlistener.com/opinion/117927/united-states-v-lopez/}{CourtListener}).

\begin{quote}
``To uphold the Government's contentions here, we would have to pile inference upon inference in a manner that would bid fair to convert congressional authority under the Commerce Clause to a general police power of the sort retained by the States.''
\end{quote}
\noindent --- Chief Justice Rehnquist, Opinion of the Court (\href{https://www.law.cornell.edu/supct/html/93-1260.ZO.html}{Cornell LII}; \href{https://www.courtlistener.com/opinion/117927/united-states-v-lopez/}{CourtListener}).

\begin{quote}
``Were the Federal Government to take over the regulation of entire areas of traditional state concern, areas having nothing to do with the regulation of commercial activities, the boundaries between the spheres of federal and state authority would blur and political responsibility would become illusory.''
\end{quote}
\noindent --- Justice Kennedy, Concurring, joined by Justice O'Connor (\href{https://www.law.cornell.edu/supct/html/93-1260.ZC.html}{Cornell LII}).

\begin{quote}
``Put simply, much if not all of Art. I, \S 8 (including portions of the Commerce Clause itself) would be surplusage if Congress had been given authority over matters that substantially affect interstate commerce. An interpretation of cl. 3 that makes the rest of \S 8 superfluous simply cannot be correct.''
\end{quote}
\noindent --- Justice Thomas, Concurring (\href{https://www.law.cornell.edu/supct/html/93-1260.ZC1.html}{Cornell LII}).

\begin{quote}
``Upholding this legislation would do no more than simply recognize that Congress had a `rational basis' for finding a significant connection between guns in or near schools and (through their effect on education) the interstate and foreign commerce they threaten.''
\end{quote}
\noindent --- Justice Breyer, Dissenting, joined by Justices Stevens, Souter, and Ginsburg (\href{https://www.law.cornell.edu/supct/html/93-1260.ZD2.html}{Cornell LII}).


\end{document}
